
\subsection{Вычисление критериев уязвимости береговой линии}

Результатом вычислений описанных в разделе~\ref{eq:CWEF} будет являться полная статистика по экспозиции и волновой нагрузке, приходящейся на береговую линию за весь период рассматриваемых метеорологических наблюдений.
Другими словами, мы получим знания о том, под каким направлением, какой высоты и с каким периодом волна подходила к берегу, сколько ватт энергии она передала ему.
Все это мы будем знать для каждой записи метеорологических данных, каждодневных или ежечасных.
Отметим с этим очевидный факт, что чем меньше будет интервал между записями, тем качественнее рассчитанные значения будут описывать происходящие в реальности процессы.

В разделе~\ref{sec:whether_data} будут подробно описаны частные вопросы получения и доступности метеорологических данных, однако уже сейчас можно определить в первом приближении их ожидаемое количество.
Так, наиболее удобный для работы сервис Open-Meteo обеспечивает доступ к историческим метеоданым начиная с 1 января 1940 года.
Несложно посчитать, что даже с суточным интервалом количество рассчитываемых значений для каждой точке будет превышать 31.5 тысячу значений.
Само хранение такого количества записей применительно к множеству точек создаст естественные трудности в их хранении и анализе.

Естественным решением этой проблемы может являться расчет специальных индексов, интегрирующих собранную информацию и позволяющих ее обобщить.
Ограничиваясь, в первом приближении, необходимостью анализа волнового воздействия назовем этот показатель как \enquote{индекс волновой экспозиции}.

\subsubsection{Индекс волновой экспозиции}
\label{sec:wer_calculation}

Вычисленный для каждого расчетного пункта временной ряд CWEF содержит достаточно информации для сравнительной оценки степени механического воздействия волн на прибрежные объекты.
Однако динамика такого влияния не может быть качественно описана в отрыве от временных и частотных характеристик полученного ряда распределения CWEF.

В рамках работы предлагается схема ранжирования, основанная на четырех взаимодополняющих критериях, каждый из которых описывает отдельный аспект волновой нагрузки: 
\begin{itemize}
    \item средний уровень энергии;
    \item экстремальную (штормовую) составляющую;
    \item направленную концентрацию волновой экспозиции;
    \item межгодовую изменчивость.
\end{itemize}

Первым из перечисленных показателей служит математическое ожидание $M(\text{CWEF})$, определяющее типичный уровень волновой нагрузки на берег и вычисляемое как среднее значение по всем записям CWEF.
С учетом большого объема массива обрабатываемого  данных, среднее значение  $\overline{\text{CWEF}}$ будет обладать максимальной стабильностью, описывая общую интенсивность фонового волнового климата.
С учетом того, что собираемые сведения о ветре являются общими для некоторого набора анализируемых, локальных, точек вариация среднего значения CWEF будет определяться углом береговой линии к доминирующим направлениям ветра и \enquote{закрытостью} рассматриваемой точки к открытому морю, позволяя выделить участки берега с наиболее \enquote{тихой} и спокойной водой:
\begin{equation}\label{eq:M_CWEF}
    \overline{\text{CWEF}} = \frac{\sum_{i=1}^n{\text{CWEF}_i}}{n}.
\end{equation}

В то же время береговая эрозия в значительной мере определяется не средним, а экстремальным волновым воздействием, накопленным в период штормов. 
Объекты с высоким показателем суммарной штормовой энергии $E_\text{storm}$ при умеренном среднем CWEF -- потенциально более уязвимы к эпизодической эрозии, чем объекты с постоянно высоким, но равномерным потоком~\cite{Harley2017}.

В работе~\cite{Harley2017} вводится "integrated storm wave energy flux" как интеграл CWEF за период, пока высота волн $H_s$ превышает штормовой порог, что может быть математически выражено как:
\begin{equation}
    E_\text{storm} = \int_{H_s > H_\text{thresh}} \text{CWEF}(t)\,dt \approx \sum_{i:\,\text{CWEF}_i > H_\text{thresh}} \text{CWEF}_i \cdot \Delta t,
\end{equation}
где $H_\text{thresh}$ -- порог штормового волнения для конкретного региона, $\Delta t$ -- суточный шаг (86 400 с).

Порог штормового волнения $H_\text{thresh}$, может задаваться как отдельный гиперпараметр, определяющий предел прочности береговой линии или прибрежных инженерных конструкций.
В то же время, для относительной оценки участков берега подверженных сильному влиянию штормовой нагрузки внутри рассматриваемой зоны между собой, в научной практике~\cite{Harley2017}, допустим выбор в качестве значения $H_\text{thresh}$ 90-го перцентиля по ряду распределения CWEF.

Приведенные показатели среднего и экстремального волнового воздействия базируются исключительно на значениях CWEF, а значит учитывают только компоненту энергии по направлению нормали относительно береговой линии, что не позволяет оценить несимметричность направлений волн.
Таким образом, чувствительность берегового профиля к волновому воздействию зависит не только от интенсивности CWEF, но и от степени его направленной концентрации. 

Точки, для которых подавляющая доля волновой энергии поступает из одного узкого сектора направлений, более уязвимы к вдольбереговой абразии, эрозионному отклику и общей перестройке берега при смене господствующих ветров или аномальных штормах~\cite{Harley2017, Basnayake2026GlobalCoastalVulnerability}. 

Коэффициент направленной концентрации волновой экспозиции $K_\text{dir}$ может быть получен как доля суммарного CWEF, приходящаяся на три доминирующих смежных румба (получаемых группировкой значений по розе ветров), относительно суммарного CWEF по всем направлениям:
\begin{equation}
    K_\text{dir} = \frac{\sum_{j \in \text{top-3}} \text{CWEF}_j^{\Sigma}}{\sum_{\text{all}} \text{CWEF}_j^{\Sigma}}.
\end{equation}

Значение $K_\text{dir}$ близкое к 1 соответствует моноориентированной экспозиции ветровой нагрузки, а близкое к 0 -- равномерному «круговому» волновому воздействию.

Долгосрочная устойчивость берегового объекта определяется также стабильностью волновой нагрузки во времени.
В модели глобального CVI~\cite{Guannel2022USVI_CVI, Basnayake2026GlobalCVI} подчеркивается, что береговые объекты, подверженные значительной межгодовой вариабельности CWEF, более чувствительны к климатическим аномалиям и трендам~\cite{Ghanavati2023_wave_surge_coastlines}. 

Изменчивость признака CWEF можно оценить через безразмерный коэффициент вариации:
\begin{equation}
    \text{CV} = \frac{\sigma_\text{CWEF}}{\overline{\text{CWEF}}},
\end{equation}
где $\sigma_\text{CWEF}$ -- стандартное отклонение суточного ряда CWEF, $\overline{\text{CWEF}}$ -- математическое ожидание CWEF, рассчитанное по формуле~\eqref{eq:M_CWEF}.

Для того, чтобы комплексно оценить воздействие волн на берег, необходимо учесть описанные выше 4 показателя:
\begin{equation}
    R = f\!\left(\text{CWEF}_\text{mean},\; E_\text{storm},\; \text{CV}_\text{CWEF},\; E_\text{int}\right).
\end{equation}

Но как это сделать, если каждый из них имеет собственную размерность и диапазон принимаемых значений?

Общая логика, принятая в современных методиках оценки прибрежной уязвимости CVI, состоит в том, что вычисленные показатели разбивается на ранговые классы (обычно 5 уровней -- от «очень низкий» до «очень высокий»).
При этом границы назначаемых классов определяются либо через квантили распределения показателей совокупности всех рассматриваемых объектов, либо через конкретные пороги, определенные по результатам эмпирических наблюдений за эрозионными процессами в рассматриваемой местности.

Так как перцентильный подход обладает универсальностью и автоматически адаптируется к региональному волновому климату, в работе используется система ранжирования, представленная в таблице~\ref{tab:cwef_percentiles}.
\begin{table}
\caption{-- Ранжирование уязвимости по перцентилям CWEF}
\label{tab:cwef_percentiles}
\centering
\begin{tabular}{c l l}
\toprule
\textbf{Ранг} & \textbf{Уровень уязвимости} & \textbf{Критерий по CWEF} \\
\midrule
1 & Очень низкий  & \makecell[l]{$< 20$-й перцентиль \\ по совокупности объектов} \\
2 & Низкий        & $20$--$40$-й перцентиль \\
3 & Умеренный     & $40$--$60$-й перцентиль \\
4 & Высокий       & $60$--$80$-й перцентиль \\
5 & Очень высокий & $> 80$-й перцентиль \\
\bottomrule
\end{tabular}
\end{table}

Итоговый индекс ранга волновой экспозиции WER вычисляется как геометрическое среднее четырех частных рангов:

\begin{equation}
    \label{eq:wer_index}
    \text{WER} = \sqrt[^4]{R_1 \cdot R_2 \cdot R_3 \cdot R_4},
\end{equation}
где:
\begin{itemize}
    \item $\text{WER}$ -- Wave Exposure Rank, безразмерный индекс волновой экспозиции прибрежного объекта, принимающий значения в диапазоне $[1,\,5]$;
    \item $R_1$ -- ранг по среднегодовому значению $\overline{\text{CWEF}}$ (Вт/м), характеризующий фоновый уровень волновой нагрузки на нормаль к берегу;
    \item $R_2$ -- ранг по суммарной штормовой энергии $E_\text{storm}$ (Дж/м), характеризующий экстремальную составляющую волновой экспозиции;
    \item $R_3$ -- ранг по коэффициенту направленной концентрации $K_\text{dir}$, характеризующий степень моноориентированности волнового потока;
    \item $R_4$ -- ранг по коэффициенту вариации $\text{CV} = \sigma_\text{CWEF}/\overline{\text{CWEF}}$, характеризующий межгодовую изменчивость волновой нагрузки.
\end{itemize}

Каждый из частных рангов $R_i$ принимает целые значения от 1 до 5 в соответствии с положением объекта в перцентильном распределении по совокупности всех рассматриваемых пунктов (таблица~\ref{tab:cwef_percentiles}). 

Использование геометрического среднего в выражении~\eqref{eq:wer_index} обеспечивает симметричное взаимодействие введенных критериев.
Ни один из них не может компенсировать экстремально низкое значение другого, что предотвращает систематическое занижение итогового ранга при доминировании одного фактора над остальными~\cite{Guannel2022USVI_CVI}.

Перцентильный подход к ранжированию делает методику применимой как для замкнутых акваторий с умеренным волнением, так и для открытых прибрежных зон. 
При наличии верифицированных данных о наблюдаемой скорости береговой эрозии перцентильные границы могут быть уточнены путем регрессионного сопоставления WER с измеренными темпами отступания берега.

\subsubsection{Способы использования}
\label{sec:wer_usage}

Введенный индекс волновой экспозиции WER может быть непосредственно использован при зонировании береговой полосы, оценки приоритизации берегозащитных мероприятий, принятии разрешительных решений на строительство и реконструкцию в береговой полосе. 

Основная идея состоит в том, чтобы связать режим использования территории и требования к проектной документации с количественно оцениваемой волновой нагрузкой, а не с фиксированными нормативными буферами, одинаковыми для всех участков береговой линии.

Наиболее прямое применение -- установление дифференцированных охранных зон в зависимости от класса WER. 
В отличие от фиксированных нормативных полос (например, 20 м от уреза воды по всему берегу), WER-зонирование позволяет сделать охранную зону адекватной реальной волновой нагрузке.
От минимальных зон для участков береговой линии где WER устойчиво мал и возрастающий в зонах его нестабильных или больших значений. 
Следует отметить, что такой подход уже реализован в ряде заграничных региональных CVI-методик, где класс волновой экспозиции напрямую определяет ширину запретной зоны для строительства~\cite{Guannel2022USVI_CVI}.

При ограниченном бюджете на берегозащитные мероприятия WER дает объективный инструмент ранжирования объектов по срочности вмешательства. Например, участки берега, где WER ≥ 4, при наличии существующей инфраструктуры образуют первоочередной перечень для проектирования берегоукрепления и пляженасыщения. 
Участки с WER 2–3 могут быть отнесены к программам превентивного мониторинга без немедленного вмешательства. 
Такая логика совместима с международными стандартами управления береговой зоной (ICZM -- Integrated Coastal Zone Management) и отдельными региональными регламентами EU Floods Directive~\cite{Ghanavati2023_wave_surge_coastlines}.

Дополнительно WER может быть встроен в разрешительные процедуры на строительство, реконструкцию и изменение назначения земельных участков в береговой зоне. 

Одним из вазможных вариантов которых может являться следующий порядок~\cite{Pang2023CoastalErosionClimateChange}:
\begin{itemize}
    \item \textbf{WER 1–2}: разрешение выдается в стандартном порядке с уведомлением о волновой экспозиции.
    \item \textbf{WER 3:} обязательная оценка воздействия на береговую зону (ОВБЗ) как часть проектной документации.
    \item \textbf{WER 4:} ОВБЗ + независимая экспертиза устойчивости сооружения к экстремальным волновым нагрузкам + требование к конструктивным решениям (подъем фундамента, мобильность).
    \item \textbf{WER 5:} мораторий на новое капитальное строительство; допускается только временная инфраструктура с обязательным планом демонтажа.
\end{itemize}

Поскольку WER вычисляется из временного ряда волнового климата, он естественно поддается периодическому обновлению по мере накопления новых данных -- например, раз в 5–10 лет или после аномальных штормовых сезонов.
Это дает регламентам «встроенный» механизм адаптации к изменению климата: при устойчивом росте WER для конкретного участка регламентные требования автоматически ужесточаются, что решает известную проблему «устаревания» нормативных зон при сдвиге волнового климата~\cite{Ghanavati2023_wave_surge_coastlines}.

В российской практике WER может быть интегрирован в систему прибрежных водоохранных зон (ст. 65 ВК РФ~\cite{WaterCodeRF}), дополняя фиксированные нормативные полосы динамически вычисляемым показателем волновой нагрузки. На международном уровне он совместим со структурой показателей CVI, используемых USGS и EEA, что облегчает включение результатов в сравнительные межрегиональные оценки~\cite{Basnayake2026GlobalCVI, Guannel2022USVI_CVI}.

\subsection{Выводы по разделу}

Во разделе сформирована методическая основа оценки волнового воздействия на прибрежные объекты и предложена система показателей, пригодная для интеграции в индексы уязвимости побережья.
На основе классической модели SMB и современных представлений о спектральной структуре ветрового волнения были получены зависимости для расчета значительной высоты волны $H_s$, периода пика $T_p$ и связанных с ними характеристик волновой энергии.
Для описания энергетического воздействия введен интегральный показатель потока волновой энергии $\mathrm{CWEF}$ и его производные (штормовой интеграл $E_{\text{storm}}$, коэффициент направленности $K_{\text{dir}}$, коэффициент вариации $\mathrm{CV}_{\mathrm{CWEF}}$, интегрированный поток $E_{\text{int}}$), что позволяет количественно сопоставлять различные участки береговой линии по уровню волновой экспозиции.

Показано, что использование перцентильного ранжирования по совокупности объектов позволяет перейти от непрерывных значений $\mathrm{CWEF}$ к дискретным классам волнового воздействия.
Введенный ранговый показатель $\mathrm{WER}$ обеспечивает компактное представление совокупного влияния нескольких энергетических критериев и может быть непосредственно использован в составе комплексных индексов уязвимости (CVI) и систем поддержки принятия решений в рамках интегрированного управления прибрежной зоной.
Предложенная система критериев адаптируема к различным временным масштабам (годовые, сезонные и штормовые значения), а также к различному качеству исходных данных (реанализационные ряды ERA5, численные волновые модели, натурные наблюдения).

Таким образом, раздел задает единый энергетический каркас для последующего анализа пространственного распределения волновой экспозиции и ее влияния на прибрежную инфраструктуру.
Разработанные показатели и схема ранжирования служат связующим звеном между физикой ветрового волнения и практическими задачами оценки риска, что обеспечивает воспроизводимость расчетов и возможность расширения методики на новые регионы и наборы данных.
